\documentclass[11pt,letterpaper]{article}

\usepackage[spanish,es-nosectiondot]{babel}
\usepackage[utf8]{inputenc}
\usepackage[T1]{fontenc}
\usepackage{amsmath}
\usepackage{booktabs}
\usepackage{geometry}
\usepackage{longtable}
\usepackage{pdflscape}
\usepackage{array}
\usepackage{tabularx}
\usepackage{threeparttable}
\usepackage{xurl}

\geometry{margin=2.5cm}

\title{Avance frente econométrico: modelo MNL con zona de residencia}
\author{Juan Vicente Onetto Romero}
\date{Mayo 2026}

\begin{document}

\maketitle

\section{Cambio metodológico}

Se actualizó la forma de medir las variables territoriales del modelo:

\begin{itemize}
    \item Las variables que representan características del usuario/hogar se miden en la \textbf{zona de residencia inferida del usuario}.
    \item Las variables que representan el contexto del viaje se mantienen en la \textbf{zona de origen del viaje}.
    \item La residencia se infiere desde viajes declarados como \texttt{HOGAR}, usando la zona de destino más frecuente por tarjeta.
    \item Para la estimación principal se usan solo tarjetas donde esa residencia es poco ambigua: casos con una única zona observada como hogar, o con una zona claramente dominante entre los viajes declarados como \texttt{HOGAR}.
\end{itemize}

Esta separación evita atribuir al usuario características de una zona donde solamente inició un viaje.

\section{Validación de la residencia inferida}

Se validó la residencia con un criterio independiente: para cada tarjeta se revisó si la zona de origen más frecuente en la mañana coincide con la zona de destino más frecuente en la tarde. Cuando ambas zonas coinciden, esa zona se interpreta como una señal alternativa de residencia.

\begin{table}[htbp]
\centering
\caption{Validación de la residencia inferida}
\small
\setlength{\tabcolsep}{4pt}
\begin{tabularx}{\textwidth}{>{\raggedright\arraybackslash}X r r r}
\toprule
Criterio & Tarjetas & Coinciden & Acuerdo \\
\midrule
Origen frecuente mañana = destino frecuente tarde & 1,130,241 & 1,086,898 & 96.2\% \\
\bottomrule
\end{tabularx}
\end{table}

Lectura: cuando el criterio alternativo entrega una señal clara, coincide en 96.2\% de los casos con la residencia inferida desde los viajes \texttt{HOGAR}. Esto respalda usar esa residencia como base para medir variables socioeconómicas.

\section{Comparación de ajuste}

Se compara el modelo candidato por residencia contra el baseline previo, que medía las mismas variables socioeconómicas por zona de origen del viaje. La comparación usa la misma muestra.

\begin{table}[htbp]
\centering
\caption{Comparación de ajuste}
\small
\setlength{\tabcolsep}{4pt}
\begin{tabularx}{\textwidth}{l >{\raggedright\arraybackslash}X r r r l}
\toprule
Modelo & Medición socioeconómica & N viajes & LL final & AIC & Estado \\
\midrule
Baseline previo & Zona de origen del viaje & 279,643 & -137,382.9 & 274,897.8 & Converge \\
Candidato main & Zona de residencia inferida & 279,643 & -136,351.6 & 272,835.2 & Converge \\
\bottomrule
\end{tabularx}
\end{table}

Lectura: medir las variables socioeconómicas por residencia mejora el ajuste del MNL en la misma muestra.

\section{Coeficientes seleccionados del modelo candidato}

La alternativa base normalizada es BIP. La tabla muestra variables seleccionadas por relevancia sustantiva; la tabla completa de parámetros queda en el anexo.

\begin{table}[htbp]
\centering
\caption{Coeficientes seleccionados del modelo candidato}
\small
\setlength{\tabcolsep}{5pt}
\begin{tabular}{p{5.4cm}rrrr}
\toprule
Variable & QR\_OTHER Est. & QR\_OTHER t & QR\_RED Est. & QR\_RED t \\
\midrule
Educación univ. o más en residencia & 0.158 & 13.34 & 0.397 & 16.88 \\
Discapacidad en residencia & 0.073 & 5.47 & -0.243 & -7.98 \\
Inmigrantes en residencia & 0.060 & 7.75 & 0.006 & 0.35 \\
Edad promedio en residencia & -0.100 & -10.27 & 0.040 & 1.87 \\
Proxy ingreso E-D en residencia & 0.007 & 0.70 & 0.041 & 1.90 \\
Macrozona Oriente & -0.004 & -0.18 & 0.122 & 2.75 \\
Tiempo de espera inicial & -0.00127 & -15.56 & -0.00104 & -10.53 \\
Tiempo espera transbordo & -0.00065 & -5.61 & -0.00029 & -2.45 \\
\bottomrule
\end{tabular}
\end{table}

Lectura breve de la tabla:
\begin{itemize}
    \item La educación universitaria residencial es la señal socioeconómica más fuerte, especialmente para QR\_RED.
    \item Macrozona Oriente muestra asociación positiva con QR\_RED y no muestra señal robusta para QR\_OTHER.
    \item El proxy de ingreso E-D no aparece como efecto directo fuerte; probablemente parte de su señal queda absorbida por educación, macrozona y otros controles.
    \item Discapacidad residencial aparece como una señal importante y no trivial; conviene interpretarla con cautela.
\end{itemize}

\section{Lectura principal}

\begin{itemize}
    \item La residencia inferida es consistente con una validación independiente basada en patrones mañana/tarde.
    \item Medir variables socioeconómicas por residencia mejora el ajuste respecto del baseline por origen.
    \item La educación residencial es el resultado más estable: zonas con mayor educación universitaria se asocian con mayor uso de QR, especialmente QR\_RED.
    \item La macrozona Oriente muestra una señal positiva para QR\_RED una vez que las características socioeconómicas se miden por residencia y no por el punto de inicio del viaje.
    \item La variable de discapacidad aparece como una señal importante y no trivial; conviene discutirla con cautela y contrastarla con el frente ML.
\end{itemize}

\section{Fuentes de resultados}

\begin{itemize}
    \item Modelo candidato por residencia: resultados Biogeme del notebook 17, especificación candidata por residencia.
    \item Baseline comparable por origen: resultados Biogeme del notebook 17, misma muestra y variables medidas por origen.
    \item Validación AM/PM: auditoría \texttt{commute\_vs\_proposito\_summary\_ml\_2025.csv}.
\end{itemize}

\appendix
\section{Anexo: parámetros completos del modelo candidato en formato wide}

Los coeficientes de variables comunes al viaje o a la zona se interpretan respecto de BIP. Para variables alternativas-específicas, como tiempos y transbordos, también se estima directamente un coeficiente para BIP.

\begin{landscape}
\scriptsize
\setlength{\tabcolsep}{2pt}
\begin{longtable}{p{0.25\linewidth}rrrrrr}
\caption{Parámetros completos del modelo candidato en formato wide} \\
\toprule
Variable & BIP Est. & BIP t & QR\_OTHER Est. & QR\_OTHER t & QR\_RED Est. & QR\_RED t \\
\midrule
\endfirsthead
\toprule
Variable & BIP Est. & BIP t & QR\_OTHER Est. & QR\_OTHER t & QR\_RED Est. & QR\_RED t \\
\midrule
\endhead
\bottomrule
\endfoot
Constante específica alternativa & & & -2.048 & -18.7294 & -4.2339 & -18.8435 \\
Tiempo en vehículo & -5.81e-05 & -1.4005 & 2.76e-05 & 0.672 & 1.9e-05 & 0.4561 \\
Espera inicial & -0.0012 & -14.8115 & -0.0013 & -15.5553 & -0.001 & -10.5288 \\
Espera transbordo & -0.0006 & -5.0754 & -0.0006 & -5.6091 & -0.0003 & -2.4549 \\
N transbordos & 0.0014 & 0.0253 & -0.0614 & -1.1468 & -0.1037 & -1.8624 \\
Laboral PM & & & 0.0814 & 4.3442 & 0.4521 & 11.4155 \\
Laboral punta tarde & & & 0.2008 & 8.9388 & 0.4106 & 9.2001 \\
No laboral & & & 0.1786 & 7.7629 & 0.1726 & 3.4562 \\
Año 2025 & & & 0.3432 & 29.5267 & 0.2802 & 11.6735 \\
Demanda zona-franja & & & -0.0341 & -3.3534 & 0.0217 & 1.0647 \\
Densidad paraderos bus & & & 0.0181 & 1.0612 & -0.1306 & -3.4147 \\
Líneas bus & & & 0.0289 & 1.574 & 0.0597 & 1.5668 \\
Líneas metro & & & 0.1109 & 4.0922 & -0.0808 & -1.5423 \\
Residencia: educación universitaria o más & & & 0.1583 & 13.3413 & 0.3973 & 16.8769 \\
Residencia: discapacidad & & & 0.0727 & 5.4689 & -0.2431 & -7.975 \\
Residencia: inmigrantes & & & 0.0598 & 7.7512 & 0.0063 & 0.3462 \\
Residencia: mujeres & & & 0.0014 & 0.269 & 0.0165 & 1.1363 \\
Residencia: asistencia parvularia & & & -0.038 & -4.143 & 0.0541 & 2.865 \\
Residencia: edad promedio & & & -0.0996 & -10.2716 & 0.0402 & 1.8749 \\
Residencia: proxy ingreso E-D & & & 0.0068 & 0.7011 & 0.0408 & 1.9038 \\
OSM: juegos/plazas infantiles & & & 0.014 & 2.2994 & 0.0423 & 3.4265 \\
OSM: colegios & & & 0.0109 & 1.6523 & -0.0172 & -1.1761 \\
OSM: universidades & & & -0.0216 & -8.064 & -0.012 & -2.1181 \\
OSM: refugios/paraderos & & & -0.0308 & -3.3724 & 0.0316 & 1.5739 \\
OSM: accesos metro & & & -0.0159 & -3.7541 & 0.0054 & 0.6346 \\
Macrozona norte & & & -0.0308 & -0.9753 & -0.0317 & -0.4097 \\
Macrozona poniente & & & -0.0591 & -2.6301 & -0.2849 & -5.3291 \\
Macrozona oriente & & & -0.0039 & -0.1783 & 0.1223 & 2.754 \\
Macrozona sur & & & -0.0474 & -1.6773 & -0.0783 & -1.1975 \\
Macrozona suroriente & & & -0.0752 & -2.9964 & -0.05 & -0.9171 \\
Macrozona externa/especial & & & 0.0275 & 0.2255 & 0.2998 & 1.3577 \\
\end{longtable}
\end{landscape}

\end{document}
